
\subsubsection{2.1 Alignment Evaluation Benchmarks}

Current alignment evaluation infrastructure focuses predominantly on AI-to-human alignment. RewardBench \citep{Lambert2024} evaluates reward models on preference prediction across chat, safety, and reasoning domains. AlignBench \citep{AlignBench2024} provides Chinese-language alignment evaluation with fine-grained capability assessment. PERSONA (2024) introduces pluralistic alignment by testing whether models can serve diverse user profiles. These benchmarks share a common assumption: alignment success means the AI accurately predicts and satisfies human preferences.

What these benchmarks do not measure is whether AI responses support or undermine human agency.

\subsubsection{2.2 Human Agency in AI Systems}

The concept of agency-preserving AI has emerged from both AI safety and human-computer interaction research. Mitelut et al. (2023) provide a framework for understanding how intent-aligned AI systems can deplete human agency by removing decision-making friction that serves epistemic purposes.

HumanAgencyBench \citep{HumanAgencyBench2024} operationalizes this concern with six measurable dimensions: clarifying ambiguity, presenting options, epistemic hedging, explicit deferral, supporting user reasoning, and avoiding cognitive shortcuts. Four of these dimensions are adapted as computable proxies in this work.

\subsubsection{2.3 Bidirectional Alignment Framework}

Shen et al. (2024) synthesize concerns about unidirectional alignment into a Bidirectional Alignment Framework. Their systematic review of 400+ papers demonstrates that the field overwhelmingly focuses on $AI\rightarrow Human$ alignment while neglecting $Human\rightarrow AI$ alignment dimensions. This work directly addresses the operationalization gap.

\subsubsection{2.4 Adversarial Probing}

Gradient reversal training, introduced by Ganin and Lempitsky (2015) for domain adaptation, provides a method for isolating orthogonal representational components. Gradient reversal is applied here to investigate whether BAI occupies a representationally independent subspace from reward prediction.
